%% kapa: exam=1 round=34 year=2023 period=2 form=A subject=appraisal_law subject_ko=감정평가관계법규 seq=24 number=24
%% source: Q-Net 감정평가사 기출문제 (cst00309) artlSeq=5211789 file=raw/1008/5211789/2023년도 제34회 감정평가사 1차시험 2교시.pdf
%% conversion: kordoc 4.24.51 -> tools/gen1.py (auto); review=proofread-sonnet; flags=none
\begin{question}{24}
건축법령상 이행강제금에 관한 설명으로 옳은 것을 모두 고른 것은?
\begin{framedtext}
ㄱ. 허가권자는 시정명령을 받은 자가 이를 이행하면 새로운 이행강제금의 부과를 즉시 중지하되, 이미 부과된 이행강제금은 징수하여야 한다. \par
ㄴ. 동일인이 「건축법」에 따른 명령을 최근 2년 내에 2회 위반한 경우 부과될 금액을 100분의 150의 범위에서 가중하여야 한다. \par
ㄷ. 허가권자는 최초의 시정명령이 있었던 날을 기준으로 하여 1년에 최대 3회 이내의 범위에서 그 시정명령이 이행될 때까지 반복하여 이행강제금을 부과·징수할 수 있다. \par
\end{framedtext}
\begin{choices}
  \item ㄱ
  \item ㄴ
  \item ㄱ, ㄴ
  \item ㄴ, ㄷ
  \item ㄱ, ㄴ, ㄷ
\end{choices}
\end{question}
\answer{1}
